\section{The genuine homogeneous initial ideal has the actual original Hilbert--Samuel function}

Work in the existing $S=\mathbb C[[X_1,\ldots,X_n]]$ with its genuine
variable ideal $\mathfrak n$, complex action and degree-first monomial order.
For actual nonzero $F$ let $a(F)$ be its internally constructed least
coefficient exponent, put $d(F)=|a(F)|$, and define the actual homogeneous
initial form $\operatorname{in}F=[F]_{d(F)}$, where $[F]_q$ is the genuine
homogeneous degree-$q$ component of the SAME formal series.
For an actual ideal $I$ define
\[
 \operatorname{in}I=(\operatorname{in}F:0\ne F\in I)\subset S.
\]
This is exactly the homogeneous initial ideal in
\cite[\S7, Remark~7.6 and Lemma~7.7, pp.~263--264]{bierstone-milman97-invent-math}.
We prove the actual whole diagram and Hilbert--Samuel identities internally,
without a supplied basis, rank, flat-degeneration or dimension certificate.

\begin{lemma}[True homogeneous components and literal initial exponents]
\label{hir:actual-homogeneous-initial-form-coefficients}
\uses{hir:actual-formal-initial-exponent-coefficient-product,hir:actual-formal-degreelex-initial-exponent-power-kernel}
The genuine lowest total order of actual nonzero $F$ is $d(F)$.
Its actual initial form has coefficient $[X^\alpha]F$ in degree $d(F)$
and zero in every other degree; it is nonzero and has exactly the SAME
least degree-first exponent $a(F)$.
For any true homogeneous $B$ of degree $d$ and any actual $H$, its
true degree-$q$ product component is
\[
 [HB]_q=\begin{cases}0&q<d,\\{}[H]_{q-d}B&q\ge d.\end{cases}
\]
\lean{Hironaka.HironakaActualHomogeneousInitialHS.actualHomogeneousInitialDegree,Hironaka.HironakaActualHomogeneousInitialHS.actualHomogeneousInitialDegree_eq_order,Hironaka.HironakaActualHomogeneousInitialHS.actualHomogeneousInitialForm,Hironaka.HironakaActualHomogeneousInitialHS.actualHomogeneousInitialForm_ne_zero,Hironaka.HironakaActualHomogeneousInitialHS.actualHomogeneousInitialForm_initialExponent,Hironaka.HironakaActualHomogeneousInitialHS.homogeneousComponent_mul_homogeneous_of_le,Hironaka.HironakaActualHomogeneousInitialHS.homogeneousComponent_mul_homogeneous_of_lt}
\leanok
\end{lemma}
\begin{proof}
The actual minimum exponent has a true nonzero coefficient. Every
exponent of lower total degree precedes it in the actual degree-first
order and therefore has zero coefficient. This proves that the genuine
formal total order is exactly $d(F)$. The actual homogeneous-component
coefficient formula retains the nonzero minimum coefficient and kills
all other degrees. Every exponent preceding $a(F)$ has either zero
original coefficient or zero projected coefficient, so actual minimum
uniqueness gives the SAME exponent for $\operatorname{in}F$.
For the product formula use the true finite antidiagonal convolution.
When the target exponent has total degree $q$, a nonzero coefficient
from $B$ has degree $d$, forcing the complementary exponent from $H$
to have degree $q-d$. If $q<d$ there is no such actual antidiagonal
pair. If $q\ge d$, replacing $H$ by its actual degree-$(q-d)$
component therefore leaves precisely the same degree-$q$ coefficients.
The product of the two actual homogeneous components is homogeneous
of degree $q$, so all its other coefficients vanish as well. Actual
coefficient extensionality proves the displayed full series identity.
\end{proof}

\begin{theorem}[Every actual homogeneous initial-ideal component lifts into the original whole ideal]
\label{hir:actual-homogeneous-initial-ideal-degree-lift}
\uses{hir:actual-homogeneous-initial-form-coefficients}
For EVERY actual $G\in\operatorname{in}I$ and EVERY $q\in\mathbb N$,
there exists an actual $F_q\in I\cap\mathfrak n^q$ such that
$[F_q]_q=[G]_q$.
\lean{Hironaka.HironakaActualHomogeneousInitialHS.actualHomogeneousInitialIdealIndex,Hironaka.HironakaActualHomogeneousInitialHS.actualHomogeneousInitialIdeal,Hironaka.HironakaActualHomogeneousInitialHS.actual_initial_product_degree_lift,Hironaka.HironakaActualHomogeneousInitialHS.actual_homogeneous_initial_ideal_degree_lift}
\leanok
\end{theorem}
\begin{proof}
Actual ideal-span membership internally supplies a finite original
combination $G=\sum_i H_i\operatorname{in}f_i$, with genuine nonzero
$f_i\in I$ and genuine formal coefficients $H_i$. Write $d_i=d(f_i)$.
For $q<d_i$ the preceding product formula gives a zero degree-$q$
component, and we choose zero for that summand's original lift.
For $q\ge d_i$, choose the actual original member
$[H_i]_{q-d_i}f_i\in I$. Its genuine order is at least
$(q-d_i)+d_i=q$, so it belongs to the actual whole $\mathfrak n^q$.
The true lowest-order product formula gives its degree-$q$ component
as $[H_i]_{q-d_i}\operatorname{in}f_i$, which is exactly the degree-$q$
component of $H_i\operatorname{in}f_i$ by the preceding lemma.
The finite actual sum of these chosen original lifts belongs to BOTH
$I$ and $\mathfrak n^q$. The genuine homogeneous-component map is
complex linear, so the sum's actual degree-$q$ component is exactly
$[G]_q$. This constructs the lift from actual finite original
coefficients and true product convolution; no desired lift is assumed.
\end{proof}

\begin{theorem}[Both whole genuine degree-first initial diagrams coincide]
\label{hir:actual-homogeneous-initial-ideal-diagram-equality}
\uses{hir:actual-homogeneous-initial-ideal-degree-lift,hir:actual-homogeneous-initial-form-coefficients}
The genuine actual diagrams satisfy
$\mathcal N(\operatorname{in}I)=\mathcal N(I)$.
\lean{Hironaka.HironakaActualHomogeneousInitialHS.actual_homogeneous_initial_ideal_diagram}
\leanok
\end{theorem}
\begin{proof}
If an actual exponent is the least exponent of nonzero $F\in I$,
the preceding lemma says that it is also the least exponent of the
SAME nonzero $\operatorname{in}F\in\operatorname{in}I$. This gives
one whole diagram inclusion.
Conversely let $G\in\operatorname{in}I$ be nonzero, let $a=a(G)$,
and put $q=|a|$. Apply the preceding actual degree-lift theorem to
obtain $F_q\in I\cap\mathfrak n^q$ with $[F_q]_q=[G]_q$.
The actual coefficient at $a$ survives this equality and is nonzero.
Every exponent preceding $a$ in the genuine degree-first order has
degree at most $q$. If its degree is less than $q$, its coefficient
in $F_q$ vanishes by true whole-power membership. If its degree is
$q$, its coefficient agrees with that of $G$ by the actual homogeneous
component equality, and vanishes by the true minimality of $a(G)$.
Thus $F_q$ is nonzero and has exactly least exponent $a$.
This puts the arbitrary actual exponent in $\mathcal N(I)$ and proves
the reverse whole inclusion.
\end{proof}

\begin{theorem}[Actual homogeneous initial ideals preserve every genuine Hilbert--Samuel value]
\label{hir:actual-homogeneous-initial-ideal-hs-equality}
\uses{hir:actual-homogeneous-initial-ideal-diagram-equality,hir:actual-formal-and-original-hs-standard-monomial-count}
For EVERY actual ideal $I$ and every $\ell\in\mathbb N$,
\[
 \dim_{\mathbb C}S/(I+\mathfrak n^{\ell+1})
 =\dim_{\mathbb C}S/(\operatorname{in}I+\mathfrak n^{\ell+1}).
\]
Both quotients use the genuine existing complex scalar action and
actual whole ideals. In particular the genuine homogeneous initial
ideal has exactly the actual original formal Hilbert--Samuel function.
\lean{Hironaka.HironakaActualHomogeneousInitialHS.genuineFormalHilbertSamuelValue_homogeneous_initial}
\leanok
\end{theorem}
\begin{proof}
The existing genuine formal division constructs the actual finite
standard-monomial basis of EACH whole displayed quotient. Its dimension
is therefore exactly the number of actual exponents of degree at most
$\ell$ outside that ideal's true degree-first diagram. The preceding
BOTH whole diagram inclusions make these actual finite index sets
identical. Their true finite cardinalities, and hence the existing
complex vector-space dimensions, agree. The argument covers the zero
and unit ideals and the empty coordinate family. No formal-series
scalar transport, finite rank premise or HS comparison is supplied.
\end{proof}
